% 海王星（综述）
% license CCBYSA3
% type Wiki

本文根据 CC-BY-SA 协议转载翻译自维基百科\href{https://en.wikipedia.org/wiki/Neptune}{相关文章}。

\begin{figure}[ht]
\centering
\includegraphics[width=8cm]{./figures/bfb6ab5163cda4f5.png}
\caption{1989年8月由旅行者2号拍摄的真实色彩海王星图像；\(^\text{[1]}\) 中央为大黑斑\(^\text{[a]}\)。} \label{fig_Neptun_1}
\end{figure}
海王星是绕太阳运行的第八颗也是已知最远的行星。按直径计算，它是太阳系中第四大行星，按质量计算为第三大行星，并且是密度最高的巨行星。它的质量是地球的17倍。与其邻近的冰巨行星天王星相比，海王星略小，但质量更大、密度更高。由于主要由气体和液体组成\(^\text{[21]}\)，它没有明确定义的固体表面。海王星每164.8年绕太阳运行一周，轨道距离为30.1个天文单位（45亿千米；28亿英里）。它以罗马海神命名，并拥有天文学符号♆，代表海神的三叉戟\(^\text{[e]}\)。

海王星肉眼不可见，也是太阳系中唯一一颗最初不是通过直接经验观测而发现的行星。相反，天王星轨道的异常变化使亚历克西·布瓦尔提出假说，认为其轨道受未知行星的引力扰动。在布瓦尔去世后，约翰·库奇·亚当斯和乌尔班·勒维耶分别依据布瓦尔的观测独立地数学预测了海王星的位置。随后在1846年9月23日\(^\text{[2]}\)，约翰·戈特弗里德·加勒用望远镜直接观测到了海王星，位置与勒维耶预测的位置相差不到一度。其最大卫星海卫一不久后被发现，但海王星剩余的其他卫星直到20世纪才通过望远镜观测到。

海王星与地球的距离使其视直径很小，而其距离太阳遥远则使其极为暗淡，给基于地面的望远镜研究带来挑战。直到哈勃太空望远镜以及配备自适应光学的大型地面望远镜出现，才使得详细观测成为可能。1989年8月25日飞掠海王星的旅行者2号仍然是唯一访问过这颗行星的航天器\(^\text{[22][23]}\)。与气体巨行星（木星和土星）类似，海王星的大气主要由氢和氦组成，并含有少量的碳氢化合物以及可能存在的氮，但其含有更高比例的冰，如水、氨和甲烷。与天王星类似，其内部主要由冰和岩石构成\(^\text{[24]}\)，这两颗行星通常被视为“冰巨行星”以示区分\(^\text{[25]}\)。除了瑞利散射外，外层区域中微量的甲烷使海王星呈现出淡蓝色\(^\text{[26][27]}\)。

与天王星强烈的季节性大气不同，后者可能长时间几乎没有可见特征，海王星的大气具有活跃且持续可见的天气形态。在旅行者2号1989年飞掠期间，该行星的南半球存在一个可与木星大红斑相比的大黑斑。2018年，新的主要暗斑及更小的暗斑被识别并研究\(^\text{[28]}\)。这些天气形态由太阳系所有行星中最强的持续风驱动，风速可高达 2{,}100 km/h（580 m/s；1{,}300 mph）\(^\text{[29]}\)。由于其与太阳的巨大距离，海王星外层大气是太阳系中最寒冷的地方之一，其云顶温度接近 55 K（−218 °C；−361 °F）。行星中心温度约为 5{,}400 K（5{,}100 °C；9{,}300 °F）\(^\text{[30][31]}\)。海王星具有微弱且破碎的环系统（称为“弧”），于 1984 年被发现并由旅行者2号确认\(^\text{[32]}\)。
\subsection{历史}
\subsubsection{发现}
\begin{figure}[ht]
\centering
\includegraphics[width=6cm]{./figures/b8d92d1f5c26ae9a.png}
\caption{伽利略·伽利莱可能是第一位观测到海王星的人} \label{fig_Neptun_2}
\end{figure}
\begin{figure}[ht]
\centering
\includegraphics[width=6cm]{./figures/a7ed05f7bd7b8f95.png}
\caption{乌尔班·勒维耶大体上成功预测了海王星的位置} \label{fig_Neptun_3}
\end{figure}
\begin{figure}[ht]
\centering
\includegraphics[width=6cm]{./figures/bb642f227f35d3c1.png}
\caption{约翰·戈特弗里德·加勒受勒维耶之请求寻找海王星} \label{fig_Neptun_4}
\end{figure}
\begin{figure}[ht]
\centering
\includegraphics[width=6cm]{./figures/6e8babe86270c7ba.png}
\caption{约翰·库奇·亚当斯独立计算了海王星的位置} \label{fig_Neptun_5}
\end{figure}
一些已知最早的望远镜观测记录是伽利略于1612年12月28日和1613年1月27日（新历）绘制的图，它们包含的定位点与现今已知的海王星在这些日期的位置相吻合。两次观测中，伽利略似乎都把海王星误认为一颗恒星，因为它在夜空中出现在靠近木星的位置——即形成合相\(^\text{[33]}\)。因此，他不被视为海王星的发现者。在1612年12月的首次观测中，海王星在天空中几乎保持静止，因为它就在那一天开始逆行。这种视逆行运动是当地球轨道超过外行星时产生的。由于海王星只是刚刚开始其每年的逆行周期，这颗行星的运动非常微弱，以至于伽利略的小型望远镜无法检测到\(^\text{[34]}\)。2009年的一项研究表明，伽利略至少意识到他所观测到的这颗“恒星”相对于固定恒星发生了移动\(^\text{[35]}\)。

1821年，亚历克西·布瓦尔发表了关于天王星轨道的天文表\(^\text{[36]}\)。随后观测揭示出与天文表相比存在显著偏差，这促使布瓦尔假设有未知天体通过引力相互作用扰动了天王星的轨道\(^\text{[37]}\)。1843年，约翰·库奇·亚当斯开始利用他所拥有的数据研究天王星的轨道。他向皇家天文学家乔治·艾里爵士请求额外数据，艾里于1844年2月提供了这些数据。亚当斯在1845—1846年继续工作，并提出了多个关于天王星之外未发现行星位置的不同估算\(^\text{[38][39]}\)。

与亚当斯独立地，乌尔班·勒维耶在1845—1846年自行得出了指向未发现行星的计算结果，但并未激起其同胞的热情。1846年6月，在看到勒维耶首次发表的可疑未发现行星经度估算及其与亚当斯估算的相似性后，艾里说服詹姆斯·查利斯进行搜索。查利斯在整个8月和9月徒劳地搜索天空\(^\text{[37][40]}\)。事实上，查利斯比后来真正发现该行星的约翰·戈特弗里德·加勒早一年观测到了海王星，并在1845年8月4日和12日两次观测到。然而，他过时的星图和糟糕的观测技术使他在后续分析之前未能识别这些观测结果。查利斯充满悔意，但将自己的疏忽归咎于星图以及同时进行彗星观测工作的干扰\(^\text{[41][37][42]}\)。

与此同时，勒维耶寄出了一封信，并敦促柏林天文台天文学家加勒使用天文台的折射望远镜进行搜索。海因里希·德阿雷斯特（Heinrich d'Arrest），一位在天文台学习的学生，向加勒建议，他们可以将雷维耶预测位置区域的最新绘制星图与当前星空进行比较，以寻找行星特有的位移，而非恒星那样的固定位置。1846年9月23日傍晚，也就是加勒收到信件的当天，他就在宝瓶座ι星东北方向、距离勒维耶所预测的“摩羯座δ星以东五度”位置1°处发现了海王星\(^\text{[43][44]}\)，该位置距离亚当斯的预测约12°，并根据现代国际天文学联合会（IAU）星座边界处于宝瓶座与摩羯座的交界处。
\begin{figure}[ht]
\centering
\includegraphics[width=6cm]{./figures/80ca30e615e65bcf.png}
\caption{加勒用于发现海王星的9英寸折射望远镜} \label{fig_Neptun_6}
\end{figure}
在发现之后，法国人与英国人之间围绕谁应当获得这一发现的荣誉产生了民族主义式的竞争。最终，国际共识形成，认为勒维耶和亚当斯应当共同享有荣誉\(^\text{[45]}\)。自1966年以来，丹尼斯·罗林斯（Dennis Rawlins）一直质疑亚当斯共同发现主张的可信性，而随着1998年“海王星文献”（历史文件）归还格林尼治皇家天文台，该问题又被历史学家重新评估\(^\text{[46][47]}\)。
\subsubsection{命名}
在其发现后不久，海王星被简单地称为“天王星外侧的行星”或“勒维耶的行星”。第一个命名建议来自加勒，他提议命名为“雅努斯”。在英国，查利斯提出了“欧绪安努斯”这一名称\(^\text{[48]}\)。

勒维耶宣称自己拥有为其发现命名的权利，迅速为这颗新行星提出了“海王星”这一名称，尽管他虚假声称该名称已获得法国经度局的正式批准\(^\text{[49]}\)。在十月，他试图以自己的名字“勒维耶”命名该行星，并得到了天文台台长弗朗索瓦·阿拉戈的忠实支持。这一建议在法国之外遭到了强烈抵制\(^\text{[50]}\)。法国年鉴很快重新启用了“赫歇尔”一名来指代天王星，以纪念该行星的发现者威廉·赫歇尔爵士，并用“勒维耶”来指代新行星\(^\text{[51]}\)。

斯特鲁维于1846年12月29日在圣彼得堡科学院赞成使用“海王星”这一名称\(^\text{[52]}\)，这一命名是基于通过望远镜观察到的行星颜色\(^\text{[53]}\)。很快，“海王星”成为国际公认的名称。在罗马神话中，海王星是海神，与希腊神话中的波塞冬相对应。对神话名称的需求似乎与其他行星的命名法相一致，后者均以希腊和罗马神话中的神祇命名\(^\text{[f][54]}\)。

今天，大多数语言都使用“海王星”名称的某种变体。在中文、越南语、日语和韩语中，该行星名称被翻译为“海王星”\(^\text{[55][56]}\)。在蒙古语中，海王星被称为“Dalain van”（Далайн ван），反映其同名神祇作为海之统治者的角色。在现代希腊语中，该行星被称为“Poseidon”（Ποσειδώνας, Poseidonas），即希腊对应神波塞冬\(^\text{[57]}\)。在希伯来语中，“Rahab”（רהב）来自《诗篇》中提到的圣经海怪，在2009年由希伯来语言学院管理的投票中被选为该行星的官方名称，尽管现有的拉丁语形式“Neptun”（נפטון）更常使用\(^\text{[58][59]}\)。在毛利语中，该行星被称为“Tangaroa”，取自毛利海神之名\(^\text{[60]}\)。在纳瓦特尔语中，该行星被称为“Tlāloccītlalli”，取自雨神特拉洛克之名\(^\text{[60]}\)。在泰语中，海王星被称为西化名称“Dao Nepchun/Nepjun”（ดาวเนปจูน），但也被称为“Dao Ket”（ดาวเกตุ，意为“计都之星”），取自印度占星术中扮演角色的下降交点计都（केतु）。在马来语中，“Waruna”这一名称自20世纪70年代起有使用\(^\text{[61]}\)，取自印度海神之名，但最终被拉丁语形式“Neptun”（马来西亚用法\(^\text{[62]}\)）或“Neptunus”（印度尼西亚用法\(^\text{[63]}\)）所取代。

常用的形容词形式为“Neptunian”。一度出现过“Poseidean”（/pəˈsaɪdiən/）这一临时形式，取自“Poseidon”，尽管“Poseidon”的常用形容词形式是“Poseidonian”（/ˌpɒsaɪˈdoʊniən/）\(^\text{[5][64]}\)。
\subsubsection{地位}
从1846年被发现起直到1930年冥王星被发现之前，海王星是已知的最远行星。冥王星被发现后被视为行星，因此海王星成为第二远的已知行星，唯有在1979年至1999年这20年间，由于冥王星的椭圆轨道使其比海王星更靠近太阳，使海王星在此期间成为距离太阳第九颗行星\(^\text{[65][66]}\)。对冥王星质量的估算越来越准确，从地球的十倍降至远低于月球，并且1992年柯伊伯带的发现，使许多天文学家开始讨论冥王星应被视为行星还是柯伊伯带的一部分\(^\text{[67][68][69]}\)。2006年，国际天文学联盟首次定义了“行星”一词，将冥王星重新分类为“矮行星”，使海王星再次成为太阳系中已知最外侧的行星\(^\text{[70]}\)。
\subsection{物理特征}
\begin{figure}[ht]
\centering
\includegraphics[width=8cm]{./figures/dbcd02bcd550af2d.png}
\caption{海王星与地球的尺寸比较} \label{fig_Neptun_7}
\end{figure}
海王星的质量为 \(1.024\times10^{26}\,\text{kg}\)\(^\text{[8]}\)，介于地球和更大的气体巨行星之间：其质量是地球的17.15倍，但仅为木星的1/19\(^\text{[g]}\)。在1巴压力处，其重力为 \(11.27\,\text{m/s}^2\)，是地球表面重力的1.15倍\(^\text{[8]}\)，仅次于木星\(^\text{[71]}\)。海王星的赤道半径为 \(24{,}764\,\text{km}\)\(^\text{[11]}\)，几乎是地球的四倍。海王星与天王星一样，是冰巨行星（巨行星的一个子类），因为它们比木星和土星更小，并且含有更高浓度的挥发物\(^\text{[72]}\)。在寻找系外行星的研究中，海王星常被用作转喻：那些发现的质量相近的天体通常被称为“海王星类”\(^\text{[73]}\)，就像科学家将各种系外天体称为“类木星”一样。
\subsubsection{内部结构}
海王星的内部结构类似于天王星。其大气占其质量的约5\%到10\%，并向内部延伸至距核心大约10\%到20\%的位置。大气层中的压力约达到 \(10\,\text{GPa}\)，即约 \(10^5\) 个大气压。大气层较低区域中含有更高浓度的甲烷、氨和水\(^\text{[30]}\)。
\begin{figure}[ht]
\centering
\includegraphics[width=8cm]{./figures/5b0e22d5659369d9.png}
\caption{以伪彩色呈现的海王星内部物理组成及其周围结构示意图} \label{fig_Neptun_8}
\end{figure}
地幔的质量相当于10到15个地球质量，且富含水、氨和甲烷\(^\text{[2]}\)。按照行星科学的惯例，这种混合物被称为“冰质”，尽管它实际上是一种高温、致密的超临界流体。这种具有高电导率的流体有时被称为“水—氨海”\(^\text{[74]}\)。地幔可能由一层离子水组成，其中水分子分解为氢离子和氧离子的“汤”；在更深处则存在超离子水，其中氧形成晶格，但氢离子在氧晶格内自由移动\(^\text{[75]}\)。在7{,}000公里深处，条件可能使甲烷分解成钻石晶体，并像冰雹一样向下坠落\(^\text{[76][77][78]}\)。科学家认为，这种“钻石雨”也出现在木星、土星和天王星上\(^\text{[79][77]}\)。劳伦斯利弗莫尔国家实验室的超高压实验表明，地幔顶部可能是一片由液态碳构成的“海洋”，其中漂浮着固态“钻石”\(^\text{[80][81][82]}\)。

海王星的核心很可能由铁、镍和硅酸盐组成，内部模型显示其质量约为地球的1.2倍\(^\text{[24]}\)。中心压力为7兆巴（700 GPa），约为地球中心压力的两倍，温度可能为5{,}400 K（5{,}100 °C；9{,}300 °F）\(^\text{[30][31]}\)。
\subsubsection{大气}
\begin{figure}[ht]
\centering
\includegraphics[width=8cm]{./figures/247ec18a4a8a7802.png}
\caption{1994年至2020年的海王星云层变化\(^\text{[83]}\)。该伪彩色图像基于哈勃太空望远镜WFPC2与WFC3仪器的数据。} \label{fig_Neptun_9}
\end{figure}
\begin{figure}[ht]
\centering
\includegraphics[width=8cm]{./figures/8072a9d15d9dd52f.png}
\caption{高空云带在海王星较低层的云面上投下阴影。为使云层更清晰，图像颜色被夸张处理。} \label{fig_Neptun_10}
\end{figure}
在高空，海王星大气由80\%的氢和19\%的氦组成\(^\text{[30]}\)，并含有微量甲烷。甲烷在光谱$600\,nm$以上的波长（红光和红外区）具有显著的吸收带。与天王星一样，大气甲烷对红光的吸收是造成海王星微弱蓝色调的原因之一\(^\text{[84]}\)，但由于天王星大气中浓集的薄雾，海王星的蓝色更为明显\(^\text{[85][86]}\)。

海王星大气分为两个主要区域：随高度上升温度下降的对流层，以及随高度上升温度升高的平流层。二者之间的边界，即对流层顶，位于0.1\,bar（10\,kPa）的压力处\(^\text{[25]}\)。随后平流层在低于 \(10^{-5}\) 至 \(10^{-4}\,bar\)（1至10\,\text{Pa}）的压力下过渡至热层\(^\text{[25]}\)。热层逐渐过渡到外逸层\(^\text{[87]}\)。

模型显示，海王星的对流层中存在不同成分的带状云层，其成分随高度改变\(^\text{[83]}\)。上层云位于低于1\,bar的压力处，在该处温度适合于甲烷凝结。介于1至5 bar（100至500\,kPa）之间，认为会形成氨和硫化氢云层。在超过5\,bar压力之上，云层可能由氨、硫氨、硫化氢和水组成。约50\,bar（5.0\,MPa）压力处应存在水冰云层，在此温度达到273\,K（$0\,^\circ C$；$32\,^\circ F$）。在更下方可能存在氨和硫化氢云层\(^\text{[88]}\)。

已观测到海王星的高空云层在下方不透明云层上投下阴影。存在环绕行星、纬度固定的高空云带。这些环状带宽度为50–150\,km，位于云面上方约50-110 km处\(^\text{[89]}\)。这些高度属于发生天气活动的对流层；天气并不发生在更高的平流层或热层。2023年8月，海王星高空云层消失，引发一项历时三十年的研究，该研究结合了哈勃太空望远镜和地面望远镜的观测。研究表明，海王星的高空云活动与太阳周期相关，而非与行星季节相关\(^\text{[83][90][91]}\)。

海王星光谱显示，其较低平流层由于甲烷紫外光解产物（如乙烷和乙炔）的凝结而呈薄雾状\(^\text{[25][30]}\)。平流层中存在微量一氧化碳和氰化氢\(^\text{[25][92]}\)。由于更高浓度的烃类，海王星平流层比天王星更温暖\(^\text{[25]}\)。

目前原因仍不清楚，行星热层具有约750\,K（$477\,^\circ\text{C}$；$890\,^\circ\text{F}$）的异常高温\(^\text{[93][94]}\)。距离太阳过远，无法由紫外辐射产生如此热量。一种候选加热机制是大气与行星磁场中离子的相互作用。其他候选机制包括从内部传播并在大气中耗散的重力波。热层中含有微量二氧化碳和水，这些物质可能来源于如流星体和尘埃等外部来源\(^\text{[88][92]}\)。
\subsubsection{颜色}
海王星的大气在可见光光谱中呈微弱蓝色，仅比天王星大气的蓝色稍微饱和一些。早期两颗行星的图像大大夸张了海王星的色彩对比“以更好地呈现云层、带状结构和风”，使其相比天王星的灰白色显得深蓝。由于两颗行星使用了不同成像系统进行拍摄，很难直接比较后期合成图像。随着时间推移，这一问题得以重新评估，并在2023年底进行了最全面的颜色归一化\(^\text{[86][95]}\)。
\begin{figure}[ht]
\centering
\includegraphics[width=6cm]{./figures/3441fbbef21d89a9.png}
\caption{原始2色 (橙绿色) NASA/JPL图像旅行者2号,颜色夸张[96]} \label{fig_Neptun_11}
\end{figure}
\begin{figure}[ht]
\centering
\includegraphics[width=6cm]{./figures/069fdecca82ae8b0.png}
\caption{颜色在2016年重新校准 (Justin Cowart)，保留了一些对比度增强[97]} \label{fig_Neptun_12}
\end{figure}
\begin{figure}[ht]
\centering
\includegraphics[width=6cm]{./figures/bfb6ab5163cda4f5.png}
\caption{颜色在2023年重新校准 (Patrick Irwin)，接近真实颜色[98]} \label{fig_Neptun_13}
\end{figure}
\subsubsection{磁层}
海王星的磁层由一个相对于其自转轴倾斜47°的磁场组成，并相对于行星物理中心至少偏移0.55个半径（约13{,}500\,\text{km}），这一特征与天王星的磁层相似。在旅行者2号抵达海王星之前，人们曾假设天王星侧转是导致其磁层倾斜的原因。通过比较两颗行星的磁场，科学家现在认为，这种极端取向可能是行星内部流动特征所造成的。该磁场可能由电导液体（可能是氨、甲烷和水的混合物）组成的薄球壳内的对流流体运动产生\(^\text{[88]}\)，从而引发发电机效应\(^\text{[99]}\)。

海王星磁场在磁赤道处的偶极分量约为14微特斯拉（0.14\,G）\(^\text{[100]}\)。海王星的偶极磁矩约为 \(2.2\times10^{17}\,\text{T·m}^3\)（$14\,\mu T\cdot R_{N}^{3}$，其中 \(R_{N}\) 为海王星半径）。海王星的磁场具有复杂的几何结构，其中包含来自非偶极分量的相对较大贡献，包括可能在强度上超过偶极分量的强四极分量。相比之下，地球、木星和土星的四极分量相对较小，其磁场相对于极轴的倾角也较小。海王星强四极分量可能是由磁场相对于行星中心的偏移以及磁场发电机几何约束所导致\(^\text{[101][102]}\)。

海王星的弓形激波（磁层开始减缓太阳风的区域）位于距离行星半径34.9倍处。磁层顶（磁层压力与太阳风压力平衡的区域）位于距离海王星23至26.5倍半径处。磁尾至少延伸至海王星半径的72倍，并且很可能更远\(^\text{[101]}\)。
\begin{figure}[ht]
\centering
\includegraphics[width=8cm]{./figures/2d2c8c7ed1f4bbbb.png}
\caption{左图：由哈勃太空望远镜拍摄的海王星可见光图像。右图：将哈勃图像与詹姆斯·韦布望远镜拍摄的近红外图像合成。由于极光无法在可见光波段观测，其近红外图像被渲染为青色阴影。} \label{fig_Neptun_14}
\end{figure}
旅行者2号在极紫外和无线电频率的测量显示，海王星存在微弱但复杂且独特的极光；然而，这些观测受限于时间，并未包含红外信息。随后使用哈勃太空望远镜的天文学家并未观测到极光，这与天王星更明显的极光形成对比\(^\text{[103][104]}\)。2025年3月，首次通过将哈勃太空望远镜的可见光图像与詹姆斯·韦布太空望远镜的近红外（NIR）图像合成，拍摄到了海王星的极光。相关数据取自2023年6月。詹姆斯·韦布太空望远镜尝试研究海王星大气的光谱特性，并成功探测到三氢阳离子（\(\mathrm{H_3^+}\)），该离子在极光期间生成，被视为气体巨行星和冰巨行星极光活动的明确指示物。海王星极光的性质在很大程度上受到其磁场独特性质的影响。与地球、木星或土星不同，海王星的磁极并未与行星自转极对齐，这也是其极光主要出现在中纬度区域，而非像地球或木星那样出现在极区的原因\(^\text{[105]}\)。
\subsection{气候}
\begin{figure}[ht]
\centering
\includegraphics[width=8cm]{./figures/74d4d47d4e2d9038.png}
\caption{“大黑斑”（顶部）、“滑行者”（中部白色云团）\(^\text{[106]}\)以及“小黑斑”（底部），并对比度进行了夸张处理} \label{fig_Neptun_15}
\end{figure}
海王星的天气特征是极为动态的风暴系统，风速可达近 $600\,m/s$（$2,200\,km/h$；$1{,}300\,mph$），超过超音速流\(^\text{[29]}\)。更典型地，通过追踪持续云层的运动可以发现，风速从向东的 \(20\,\text{m/s}\) 到向西的 \(325\,\text{m/s}\) 不等\(^\text{[107]}\)。在云顶，盛行风速从赤道附近的 \(400\,\text{m/s}\) 到两极的 \(250\,\text{m/s}\)\(^\text{[88]}\)。海王星上大多数风向与行星自转方向相反\(^\text{[108]}\)。总体风向模式显示，高纬度为顺行旋转，而低纬度为逆行旋转。认为这种流动方向差异是一种“表层效应”，而非更深层大气过程所致\(^\text{[25]}\)。在南纬70°，一条高速急流以 \(300\,\text{m/s}\) 的速度流动\(^\text{[25]}\)。由于季节变化，观测到海王星南半球的云带在尺寸和反照率方面有所增加。这一趋势最早在1980年被发现。海王星的轨道周期很长，其季节持续40个地球年\(^\text{[109]}\)。

海王星在典型气象活动水平上不同于天王星。旅行者2号在1989年飞掠海王星期间观测到了天气现象\(^\text{[110]}\)，但在1986年飞掠天王星时并未观测到类似现象。

在海王星赤道地区，甲烷、乙烷和乙炔的丰度比两极高10—100倍。这被解释为赤道上升运动、两极下沉运动的证据，因为仅靠光化学作用无法在缺乏经向环流的情况下解释这种分布\(^\text{[25]}\)。

2007年，人们发现海王星南极上层对流层比大气其他区域（平均约73\,K，即 $-200\,^\circ C$）高出约10\,K\(^\text{[111]}\)。这一温差足以使甲烷在靠近两极的平流层中逸出，而在大气其他区域甲烷在对流层中是冻结的。这个相对的“热点”归因于海王星的自转轴倾角，使南极在海王星过去四分之一个公转周期（约40个地球年）中一直暴露在太阳下。随着海王星缓慢移动至太阳的另一侧，南极将进入黑暗，北极将被照亮，从而使甲烷释放转移到北极\(^\text{[112]}\)。
\subsubsection{风暴}
\begin{figure}[ht]
\centering
\includegraphics[width=6cm]{./figures/9c2d66f1227d5a37.png}
\caption{旅行者2号拍摄的增强色彩图像中的“大黑斑”} \label{fig_Neptun_16}
\end{figure}
1989年，NASA 的旅行者2号航天器发现了“大黑斑”，这是一种反气旋风暴系统，范围达到 \(13{,}000\,\text{km} \times 6{,}600\,\text{km}\)（\(8{,}100\,\text{mi} \times 4{,}100\,\text{mi}\)）\(^\text{[110]}\)。该风暴类似于木星的大红斑。大约五年之后，即1994年11月2日，哈勃太空望远镜未能在行星上看到大黑斑。取而代之的是，在海王星北半球发现了一个与大黑斑相似的新风暴\(^\text{[113]}\)。

“滑行者”（Scooter）是另一种风暴，由位于大黑斑更南方的白色云团构成。这个昵称最早出现在1989年旅行者2号接近期间的几个月，当时观测到这些云团移动速度比大黑斑更快（后来获得的图像随后揭示了移动速度甚至比最初由旅行者2号探测到的云团还要快的云层）\(^\text{[108]}\)。“小黑斑”是一个南部气旋风暴，是1989年遭遇中观测到的第二强烈风暴。它最初完全呈暗色，但随着旅行者2号接近行星，其内部出现了一个明亮核心，可在大多数最高分辨率图像中看到\(^\text{[114]}\)。2018年，发现并研究了一个较新的主要暗斑和较小暗斑\(^\text{[28]}\)。2023年，首次宣布从地面观测到了海王星上的暗斑\(^\text{[115]}\)。
\begin{figure}[ht]
\centering
\includegraphics[width=10cm]{./figures/1ddab6deac2393fe.png}
\caption{} \label{fig_Neptun_17}
\end{figure}
人们认为海王星的暗斑出现在对流层中、位于明亮云层特征之下的较低高度\(^\text{[117]}\)，因此在上层云面中呈现为“空洞”。由于这些结构可以保持稳定数月之久，因而被认为是涡旋结构\(^\text{[89]}\)。暗斑常与更明亮、持久的甲烷云相关，这些云形成于对流层顶附近\(^\text{[118]}\)。伴生云的持续存在表明，即便暗斑不再呈现为可见的暗特征，一些先前的暗斑可能仍继续作为气旋存在。当暗斑向赤道迁移过近时，或通过某种未知机制，可能会消散\(^\text{[119]}\)。
\begin{figure}[ht]
\centering
\includegraphics[width=8cm]{./figures/8c867388f1228973.png}
\caption{使用 NASA/ESA 哈勃太空望远镜广域相机3拍摄的四幅图像，拍摄时间相隔数小时。近红外辐射数据被用作红色通道\(^\text{[120]}\)。} \label{fig_Neptun_18}
\end{figure}
1989年，旅行者2号的行星射电天文学（PRA）实验观测到约60次闪电，即海王星静电放电，其释放能量超过 \(7\times10^{8}J\)\(^\text{[121]}\)。等离子体波系统（PWS）在磁纬7–33°范围内探测到16次频率在50–12\,\text{kHz}之间的电磁波事件\(^\text{[122][123]}\)。这些等离子体波探测可能是在20分钟内由磁层中氨云中的闪电触发的\(^\text{[123]}\)。

在旅行者2号最近距海王星的飞掠过程中，PWS仪器以每秒28{,}800次采样率提供了海王星首次等离子体波探测\(^\text{[123]}\)。测得的等离子体密度范围为 \(10^{-3}-10^{-1}\,cm^{3}\)\(^\text{[123][124]}\)。海王星闪电可能发生在三层云中\(^\text{[125]}\)，微物理模型表明，大多数闪电可能发生在对流层的水云中，或发生在磁层浅层氨云中\(^\text{[122][126]}\)。预计海王星的闪电发生率约为木星的1/19，并且大多数闪电活动出现在高纬地区。然而，海王星闪电似乎更类似地球闪电，而非木星闪电\(^\text{[121]}\)。
\subsubsection{内部加热}
与天王星相比，海王星更为多样的天气部分归因于其更高的内部加热。海王星对流层上层的最低温度为 $51.8\,K$（$-221.3\,^\circ C$）。在大气压力等于$1\,bar$（$100\,kPa$）的深度，温度为 $72.00\,K$（$-201.15\,^\circ{C}$）\(^\text{[127]}\)。在更深的气体层中，温度稳步上升。与天王星一样，这种加热的来源仍不清楚，但差异更大：天王星仅辐射出其从太阳所接收能量的1.1倍\(^\text{[128]}\)，而海王星辐射的能量约为其从太阳接收能量的2.61倍\(^\text{[16]}\)。

海王星距离太阳比天王星远50\%以上，接收的阳光仅约为天王星的40\%\(^\text{[25]}\)；然而，其内部能量仍足以驱动太阳系中最快的行星风系。根据其内部热力特性，海王星形成时残留的热量可能足以解释其当前的热流，但在保持两颗行星明显相似的前提下，解释天王星缺乏内部热量则更加困难\(^\text{[129]}\)。
\subsection{轨道与自转}
\begin{figure}[ht]
\centering
\includegraphics[width=8cm]{./figures/a6a2c41b87bf5127.png}
\caption{海王星（红色弧线）每完成一次绕太阳（中心）轨道运行，地球需绕行164.79次。淡蓝色点代表天王星。} \label{fig_Neptun_19}
\end{figure}
海王星与太阳的平均距离为45亿千米（约30.1个天文单位（AU），即地球到太阳的平均距离），其轨道平均周期为164.79年，变动约为±0.1年。近日点距离为29.81\,AU，远日点距离为30.33\,AU\(^\text{[h]}\)。海王星轨道偏心率仅为0.008678，使其成为太阳系中轨道第二圆的行星，仅次于金星\(^\text{[131]}\)。海王星的轨道平面相对于地球轨道倾斜1.77°。

2011年7月11日，海王星完成了自1846年发现以来的首次完整质心轨道\(^\text{[132]}\)；它在天空中并未出现在精确的发现位置，因为地球在自身365.26天轨道中的位置不同。由于太阳相对于太阳系质心的运动，在7月11日，如果使用更常见的日心坐标系统，海王星相对于太阳的精确发现位置并未被到达——发现经度在2011年7月12日被到达\(^\text{[12][133][134][135]}\)。

海王星自转轴倾角为28.32°\(^\text{[136]}\)，这与地球（23°）和火星（25°）的倾角相似。因此，海王星经历与地球类似的季节变化。海王星的轨道周期很长，意味着其季节持续40个地球年\(^\text{[109]}\)。其恒星日（自转周期）约为16.11小时\(^\text{[12]}\)。由于其自转轴倾角与地球接近，其在漫长轨道周期中昼长变化并不更为极端。

由于海王星不是固体，其大气经历差异自转。宽广的赤道区域自转周期约为18小时，慢于行星磁场的16.1小时自转。相反，在极区，自转周期为12小时。该差异自转是太阳系中最显著的\(^\text{[137]}\)，并导致强烈的纬向风切变\(^\text{[89]}\)。
\subsection{形成与共振}
\subsubsection{形成}
\begin{figure}[ht]
\centering
\includegraphics[width=10cm]{./figures/bfccf4e3e1b8f2bf.png}
\caption{一项关于外行星与柯伊伯带的模拟：a）木星与土星达到2:1共振之前；b）海王星轨道改变后柯伊伯带天体向内散射之后；c）被木星逐出散射的柯伊伯带天体之后} \label{fig_Neptun_20}
\end{figure}
冰巨行星海王星和天王星的形成一直难以精确建模。当前模型表明，太阳系外部区域的物质密度过低，无法用传统认可的核吸积机制形成如此巨大的天体，因此提出了多种假设来解释它们的形成。其中一种认为冰巨行星并非由核吸积形成，而是源于原始原行星盘内部的不稳定性，随后其大气被附近大质量 OB 型恒星的辐射所剥离\(^\text{[72]}\)。

另一种观点认为，它们形成于更靠近太阳的位置，在那里物质密度更高，随后在气态原行星盘被移除之后迁移到当前轨道\(^\text{[138]}\)。这一“形成后迁移”假说因其能够更好地解释在海王星外区域观察到的小天体族群分布而更受支持\(^\text{[139]}\)。目前最广泛接受的对此假说细节的解释称为“尼斯模型”，这是一种动力学演化情景，探讨迁移中的海王星及其他巨行星对柯伊伯带结构的潜在影响\(^\text{[140][141][142]}\)。
\subsubsection{轨道共振}
\begin{figure}[ht]
\centering
\includegraphics[width=8cm]{./figures/2c58e41dd66013d3.png}
\caption{一张示意图展示了海王星在柯伊伯带中引起的主要轨道共振：高亮区域分别为2:3共振（冥卫类）、非共振“经典带”（立方卫星类），以及1:2共振（双卫类）。} \label{fig_Neptun_21}
\end{figure}
海王星的轨道对其外侧区域（即柯伊伯带）具有深远影响。柯伊伯带是一个由小型冰质天体组成的环状区域，与小行星带类似，但规模更大，从海王星轨道的30\,AU延伸到距离太阳约55\,AU\(^\text{[143]}\)。与木星引力主导小行星带相同，海王星的引力主导柯伊伯带。在太阳系形成的时间尺度上，柯伊伯带的某些区域因海王星的引力而变得不稳定，形成其结构中的空隙。40至42\,AU之间的区域就是一个例子\(^\text{[144]}\)。

在这些空隙区域内确实存在能够在太阳系寿命内保持稳定的轨道。这些轨道共振发生在海王星的轨道周期与这些天体的轨道周期呈精确分数关系时，例如1:2或3:4。比如，当某个天体每运行一圈时海王星运行两圈，那么当海王星回到原始位置时，该天体仅完成半圈。柯伊伯带中最密集的共振区域是2:3共振，已知有超过200个天体\(^\text{[145]}\)。处于该共振中的天体每完成2圈对应海王星的3圈，称为冥卫类，因为其中最大的一颗已知柯伊伯带天体冥王星属于此类\(^\text{[146]}\)。尽管冥王星经常穿越海王星轨道，但2:3共振确保它们永不相撞\(^\text{[147]}\)。3:4、3:5、4:7和2:5共振区域则较少被天体占据\(^\text{[148]}\)。

海王星存在多个已知的特洛伊天体，分别位于日–海王星L4与L5拉格朗日点——即海王星轨道前方和后方的引力稳定区域\(^\text{[149]}\)。海王星特洛伊天体可视作与海王星呈1:1共振。其中一些海王星特洛伊天体在轨道上极为稳定，可能与海王星一同形成，而非后期捕获。首个被确认位于海王星尾随L5拉格朗日点的天体是2008 LC18\(^\text{[150]}\)。海王星还有一个临时的准卫星，即（309239）2007 RW10\(^\text{[151]}\)。该天体作为海王星准卫星已持续约12{,}500年，并将在未来约12{,}500年继续保持这一动力学状态\(^\text{[151]}\)。
\subsection{卫星}
\begin{figure}[ht]
\centering
\includegraphics[width=8cm]{./figures/e47bef6571ae9463.png}
\caption{詹姆斯·韦布太空望远镜拍摄的海王星多颗卫星标注图像。明亮的蓝色衍射星状光斑是海王星最大的卫星——海卫一。} \label{fig_Neptun_22}
\end{figure}
海王星已知有16颗卫星\(^\text{[152]}\)。海卫一（Triton）是海王星最大的卫星，占据围绕海王星轨道质量的99.5\%以上\(^\text{[i]}\)，并且是唯一质量足够大而呈球形的卫星。海卫一由威廉·拉塞尔在海王星发现仅17天后发现。与太阳系中所有其他大型行星卫星不同，海卫一具有逆行轨道，这表明它是被捕获的而非原位形成；它很可能曾是柯伊伯带中的一颗矮行星\(^\text{[153]}\)。海卫一与海王星距离足够近，因此被潮汐锁定为同步自转，并由于潮汐加速而缓慢向内螺旋。约36亿年后，当其抵达洛希极限时，最终将被撕裂\(^\text{[154]}\)。1989年，海卫一是太阳系中当时已测量到的最寒冷天体\(^\text{[155]}\)，其估计温度为38\,\text{K}（−235\,^\circ\text{C}）\(^\text{[156][157]}\)。这种极低温度源于海卫一极高的反照率，使其大量反射阳光而非吸收\(^\text{[158][159]}\)。

海王星第二颗被发现的卫星（按发现顺序），即不规则卫星海卫二（Nereid），拥有太阳系中最偏心的轨道之一。其偏心率为0.7512，使其远地点距离为近地点的七倍\(^\text{[j]}\)。

从1989年7月至9月，旅行者2号发现了海王星的六颗卫星\(^\text{[160]}\)。其中，不规则形状的普罗透斯（Proteus）值得注意，因为其大小已经接近在自身引力作用下不被拉成球形天体所能达到的极限\(^\text{[161]}\)。尽管普罗透斯是海王星第二大质量的卫星，但其质量仅为海卫一的0.25\%。海王星最内侧的四颗卫星——奈阿德（Naiad）、塔拉萨（Thalassa）、得斯皮娜（Despina）和伽拉忒亚（Galatea）——运行轨道足够靠近，以至于位于海王星行星环内部。下一颗外侧卫星拉里萨（Larissa）最初于1981年通过其掩星现象被发现。该掩星现象最初被归因于行星环弧，但当旅行者2号在1989年观测海王星时，发现实际是拉里萨造成的。2002年至2003年间发现的五颗新的不规则卫星于2004年公布\(^\text{[162][163]}\)。2013年，通过合并多幅哈勃太空望远镜图像，发现了一颗新的、更小的卫星——海马（Hippocamp）\(^\text{[164]}\)。由于海王星是罗马海神，海王星的卫星均以较低级别的海神命名\(^\text{[54]}\)。
\subsection{行星环}
\begin{figure}[ht]
\centering
\includegraphics[width=8cm]{./figures/11a3797dba87348d.png}
\caption{2022年由詹姆斯·韦布太空望远镜以红外成像观测的海王星行星环与卫星。} \label{fig_Neptun_23}
\end{figure}
海王星拥有行星环系统，但其规模远小于土星和天王星的行星环\(^\text{[165]}\)。这些环可能由覆有硅酸盐或碳基物质的冰颗粒构成，这很可能使其呈现出红色调\(^\text{[166]}\)。三条主要行星环分别为距离海王星中心63{,}000\,\text{km}的狭窄亚当斯环（Adams Ring），距离53{,}000\,\text{km}的勒维耶环（Le Verrier Ring），以及距离42{,}000\,\text{km}的更宽且更暗淡的加勒环（Galle Ring）。勒维耶环向外一条微弱的延伸被命名为拉塞尔环（Lassell），其外缘由距离57{,}000\,\text{km}的阿拉戈环（Arago Ring）界定\(^\text{[167]}\)。

第一条行星环于1968年由爱德华·吉南领导的团队探测到\(^\text{[32][168]}\)。20世纪80年代初，对这一数据及更新观测的分析提出了该环可能不完整的假设\(^\text{[169]}\)。1984年一次恒星掩食中首次出现行星环可能存在间隙的证据，当时行星环在恒星进入时遮蔽了恒星，但在离开时却没有\(^\text{[170]}\)。1989年旅行者2号的图像解决了这一问题，显示了多条微弱的行星环。

最外侧的亚当斯环包含五条显著的弧段，现在命名为“勇气”“自由”“平等1”“平等2”与“博爱”（Courage, Liberté, Égalité 1, Égalité 2 和 Fraternité）\(^\text{[171]}\)。弧段的存在曾难以解释，因为力学定律预测弧段应在短时间内扩散成连续环。天文学家如今认为，弧段被内侧不远处的卫星伽拉忒亚（Galatea）的引力效应束缚于当前形态\(^\text{[172][173]}\)。

2005年宣布的地面观测结果显示，海王星行星环的稳定性远低于此前的认识。2002年和2003年从W. M. 凯克天文台拍摄的图像与旅行者2号图像相比显示行星环呈现明显衰变。特别是，自由弧（Liberté）可能在短短一个世纪内消失\(^\text{[174]}\)。
\subsection{观测}
\begin{figure}[ht]
\centering
\includegraphics[width=6cm]{./figures/ecd8d7084ae79d9d.png}
\caption{2022年海王星在宝瓶座背景恒星前的运动} \label{fig_Neptun_24}
\end{figure}
海王星在1980年至2000年期间亮度增加了约10\%，主要原因是季节变化\(^\text{[175]}\)。随着海王星在2042年接近日点，其亮度可能继续增强。其视星等目前范围为7.67至7.89，平均为7.78，标准差为0.06\(^\text{[18]}\)。在1980年之前，该行星的亮度低至8.0等\(^\text{[18]}\)。海王星过于暗淡，无法用肉眼直接观测。它的亮度甚至会被木星的伽利略卫星、矮行星谷神星以及小行星灶神星（4 Vesta）、智神星（2 Pallas）、虹神星（7 Iris）、婚神星（3 Juno）和守护神星（6 Hebe）所盖过\(^\text{[176]}\)。使用望远镜或强力双筒望远镜可以将海王星解析为一个小型蓝色圆盘，与天王星外观相似\(^\text{[177]}\)。

由于海王星与地球距离遥远，其角直径仅在2.2至2.4角秒之间变动\(^\text{[8][20]}\)，为太阳系行星中最小的。其微小的视直径使得目视观测具有挑战性。在哈勃太空望远镜和配备自适应光学（AO）的大型地基望远镜出现之前，望远镜数据非常有限\(^\text{[178][179][180]}\)。首个使用自适应光学从地基望远镜对海王星进行具有科学价值的观测始于1997年，地点为夏威夷\(^\text{[181]}\)。海王星目前正接近日点（距离太阳最近的位置），并已显示正在升温，伴随大气活动和亮度增强。随着技术进步，配备自适应光学的地基望远镜正在记录越来越多的细节图像。自20世纪90年代中期以来，哈勃望远镜和地基自适应光学望远镜在太阳系内取得了众多新发现，包括外行星周围已知卫星数量的大幅增加等。在2004年与2005年，发现了5颗新的海王星小型卫星，其直径在38至61千米之间\(^\text{[182]}\)。

从地球观测，海王星每367天经历一次视逆行，导致在每次冲日期间相对于背景恒星出现一个环形运动。这些环形轨迹曾在2010年4月与7月，以及2011年10月与11月，使海王星靠近1846年发现时的坐标\(^\text{[135]}\)。

海王星164年的轨道周期意味着该行星平均需要13年穿越黄道十二宫中的每个星座。2011年，海王星完成自发现以来首次完整绕太阳一周，并回到第一次观测位置——宝瓶座ι星东北方向\(^\text{[43]}\)。

在射电波段观测海王星显示，它同时是连续辐射和不规则爆发辐射源。这两种辐射均被认为源自其旋转磁场\(^\text{[88]}\)。在红外波段，海王星的风暴在较冷背景下显得明亮，因此可以轻松追踪这些结构的大小与形状\(^\text{[183]}\)。
\subsection{探测}
旅行者2号是唯一访问过海王星的航天器。其于1989年8月25日最接近该行星。由于这是该航天器可访问的最后一颗主要行星，任务团队决定对海王星的卫星海卫一进行一次近距离飞掠，而不考虑轨道的后果，这与旅行者1号在土星及其卫星土卫六上的遭遇操作类似。旅行者2号回传至地球的图像成为PBS在1989年播出的整夜特别节目《Neptune All Night》的基础\(^\text{[184]}\)。

在遭遇期间，航天器的信号需要246分钟才能抵达地球。因此，大多数情况下，旅行者2号的任务依赖于预先加载的海王星遭遇指令。该航天器在8月25日进入距离海王星大气4{,}400\,\text{km}范围之前，几乎与海卫二进行了近距离遭遇，随后在同一天靠近该行星最大的卫星海卫一\(^\text{[185]}\)。

航天器验证了围绕行星存在磁场，并发现其磁场偏离行星中心且在倾角上与天王星的磁场类似。通过射电辐射测量确定了海王星的自转周期，旅行者2号显示海王星具有令人惊讶的活跃天气系统。发现了6颗新卫星，并显示该行星拥有多条行星环\(^\text{[160][185]}\)。这次飞掠提供了海王星质量的首次精确测量，结果比先前计算值低0.5\%。这一新数据否定了假设存在未被发现的X行星影响海王星和天王星轨道的理论\(^\text{[186][187]}\)。

自2018年以来，中国国家航天局一直在研究一项概念，即一对类似旅行者号的星际探测器，暂称为“深梭”计划\(^\text{[188]}\)。两艘探测器将于2020年代发射，并采取不同路径探索日球层的相反方向；第二艘探测器IHP-2将于2038年1月飞掠海王星，掠过云顶高度仅1{,}000\,\text{km}，并有可能搭载在接近过程中释放的“大气撞击器”\(^\text{[189]}\)。之后，它将继续在柯伊伯带开展任务，前往目前尚未探测的日球层尾部区域。

在旅行者2号和IHP-2飞掠之后，海王星系统科学探测的下一步被认为是轨道任务；大多数提案来自NASA，多为旗舰级轨道器构想\(^\text{[190]}\)。2003年，NASA的“远景任务研究”中曾提出一项“携带探测器的海王星轨道器”任务，旨在达到卡西尼号级别的科学研究水平\(^\text{[191]}\)。随后一项未被选中的提案为“阿尔戈号”（Argo）飞掠探测器，计划于2019年发射，访问木星、土星、海王星以及一颗柯伊伯带天体。任务重点本将集中在海王星及其最大卫星海卫一，计划于2029年前后进行探测\(^\text{[192]}\)。

拟议的“新视野2号”任务可能会对海王星系统进行近距离飞掠，但后来被取消。目前Discovery计划中有一项待定提案，即“Trident”航天器，将对海王星和海卫一进行飞掠探测；\(^\text{[193]}\) 然而，该任务未被选入Discovery 15或16。另一项名为“Neptune Odyssey”的概念为海王星轨道器与大气探测器，被研究为NASA潜在的大型战略科学任务；其计划于2031年至2033年间发射，并于2049年抵达海王星\(^\text{[194]}\)。不过，出于技术规划因素，冰巨行星轨道任务中最终推荐优先级最高的是天王星轨道器与探测器，其优先级高于“土卫二轨道着陆器”（Enceladus Orbilander）\(^\text{[195]}\)。

两项值得注意的海卫一重点海王星轨道器任务提案，其成本介于Trident和Odyssey任务之间（属于“新前沿计划”），分别是“Triton Ocean World Surveyor”和“Nautilus”，其巡航阶段分别对应2031—47年和2041—56年\(^\text{[196][197]}\)。海王星也是中国“天问五号”（Tianwen-5）的潜在目标，可能于2058年前后抵达\(^\text{[198]}\)。
\subsection{另见}
\begin{itemize}
\item 海王星的轮廓
\item 热海王星
\item 海王星在占星术
\item 小说中的海王星
\item 镎
\item 海王星，神秘主义者-七个动作之一古斯塔夫·霍尔斯特的行星套房
\item 遥远未来的时间表
\item 太阳系行星的统计
\end{itemize}
\subsection{备注}\\
a.海王星的黑点不是永久的特征；旅行者2号被指定为 “1989年大黑点” 的GDS-89。\\
b.轨道元素指的是海王星重心和太阳系重心。这些是瞬时的密接精确处的值J2000纪元。给出重心量是因为与行星中心相反，它们不会因卫星的运动而每天经历明显的变化。\\
c.指1 bar (100 kPa) 大气压的水平\\
d.基于1 bar大气压水平内的体积\\
e.第二个符号，一个 'LV' 会标⯉对于 “le verrier”，类似于 “h” 会标♅为了天王星.它在法国以外从未被广泛使用，现在已经过时了。\\
t.有人可能会争辩说，这是真的，除了 “地球”，在英语中是日耳曼神的名字，Erda。的IAU政策是，人们可以用语言中常用的任何名称来称呼地球和月球。根据IAU，“terra” 和 “luna” 是不地球及其月球的官方名称:“天文物体的命名”。国际天文学联合会。已存档从2024年3月21日开始。已检索4月27日2024。\\
g.地球的质量为 \(5.9722\times10^{24}\,\text{kg}\)，由此得到质量比
\[
\frac{M_{\text{Neptune}}}{M_{\text{Earth}}}
=
\frac{1.02\times10^{26}}{5.97\times10^{24}}
=
17.15.~
\]
天王星的质量为 \(8.6810\times10^{25}\,\text{kg}\)，由此得到质量比
\[
\frac{M_{\text{Uranus}}}{M_{\text{Earth}}}
=
\frac{8.68\times10^{25}}{5.97\times10^{24}}
=
14.54.~
\]
木星的质量为 \(1.8986\times10^{27}\,\text{kg}\)，由此得到质量比
\[
\frac{M_{\text{Jupiter}}}{M_{\text{Neptune}}}
=
\frac{1.90\times10^{27}}{1.02\times10^{26}}
=
18.63.~
\]
质量数值来源：Williams, David R.（2007年11月29日），“Planetary Fact Sheet – Metric”，NASA。该文献于2014年9月5日存档，检索于2008年3月13日。\\
h.最近三次远日点距离分别为30.33\,AU，下一次为30.34\,AU。近日点距离更加稳定，为29.81\,AU\(^\text{[130]}\)。\\
i.海卫一的质量为 \(2.14\times10^{22}\,\text{kg}\)。海王星其他已知12颗卫星的合计质量为 \(7.53\times10^{19}\,\text{kg}\)，约为0.35\%。行星环的质量可以忽略不计。\\
j.\[
\frac{r_{a}}{r_{p}}
=
\frac{2}{1-e}-1
=
2/0.2488 - 1
\approx 7.039.~
\]
\subsection{参考资料}
\begin{enumerate}
\item 欧文，帕特里克G.J.;多宾森，杰克; 詹姆斯，阿朱那; 提比，尼古拉斯·A；艾米·西蒙；Fletcher, Leigh N.;罗曼，迈克尔·T；格伦·S·奥顿；Wong，Michael H.; 托莱多，丹尼尔; Pérez-hoyos，圣地亚哥; 贝克，朱莉 (2023年12月23日)。“模拟天王星的颜色和大小的季节性周期，并与海王星进行比较”。皇家天文学会月刊。527(4):11521-11538。doi:10.1093/mnras/stad3761。高密度脂蛋白:20.500.11850/657542。
\item 汉密尔顿，卡尔文J.(2001年8月4日)。"海王星"。太阳系的观点。存档自原来的2007年7月15日。已检索8月13日2007。
\item 伊丽莎白·沃尔特 (2003年4月21日)。“海王星”。剑桥高级学习词典(第二版)。剑桥大学出版社。ISBN  978-0-521-53106-1。
\item "海王星"。牛津英语词典(在线版)。牛津大学出版社。 (订阅或参与机构成员必需。)
\item “利用小型放射性同位素动力系统进行勘探”(PDF)。NASA.2004年9月。存档自原来的(PDF)2016年12月22日。已检索1月26日2016。
\item 约曼，唐纳德·K。“海王星Barycenter (Major Body = 8) 的地平线Web界面”。JPL Horizons在线星历系统。已存档从原始的2021年9月7日。已检索7月18日2014。-选择 “星历类型: 轨道元素”，“时间跨度: 2000-01-01 12:00至2000-01-02”。(“目标体: 海王星重心” 和 “中心: 太阳系重心 (@ 0)”。)
\item 塞利格曼，考特尼。“轮换周期和日长”。存档自原来的2011年7月28日。已检索8月13日2009年。
\item 威廉姆斯，大卫R.(2024年10月3日)。"海王星概况介绍"。NASA.已存档从2025年8月5日开始。已检索10月1日2025。
\item 苏阿米，D.；Souchay, J.(2012年7月)。“太阳系的不变平面”。天文学与天体物理学。543: 11。Bibcode:2012A & A...543A.133S。doi:10.1051/0004-6361/201219011。A133.
\item “地平线行星中心批次呼叫2042年9月的近日点”。ssd.jpl.nasa.gov(海王星行星中心 (899) 的近日点发生在2042年9月4日29.80647406au在rdot从负到正翻转期间)。NASA/JPL。已存档从原始的2021年9月7日。已检索9月7日2021。
\item Seidelmann，P。肯尼斯; 阿基纳，布伦特A；迈克尔·F·阿赫恩；阿尔伯特·R·康拉德；Consolmagno, Guy J.;Hestroffer，Daniel; 等人 (2007年)。“IAU/IAG制图坐标和旋转要素工作组的报告: 2006年”。天体力学与动力天文学。98(3):155-180。Bibcode:2007CeMDA .. 98 .. 155S。doi:10.1007/s10569-007-9072-y。
\item 孟塞尔，K.；史密斯，H.; 哈维，S.(2007年11月13日)。“海王星: 事实与数字”。NASA.存档自原来的2014年4月9日。已检索8月14日2007。
\item 德佩特，Imke; 利绍尔，杰克J。(2015)。行星科学(第二次更新版本)。纽约: 剑桥大学出版社。第250页。ISBN 978-0-521-85371-2。已存档从原来的2016年11月26日。已检索8月17日2016。
\item 卡罗琳·肯尼特 (2022)。天王星和海王星。Reaktion书籍。第185页。ISBN 978-1-78914-642-4。
\item 阿基纳，B。A.;阿克顿，C。H、; A'Hearn，M。F.;康拉德，A.；Consolmagno, G.J.;达克斯伯里，T.；赫斯特罗弗，D.；希尔顿，J。L.;柯克，R.L.;Klioner, S.A.;麦卡锡，D.；米奇，K.；奥伯斯特，J.；平，J.；Seidelmann，P。K。(2018)。“IAU制图坐标和旋转要素工作组的报告: 2015年”。天体力学与动力天文学。130(3): 22。Bibcode:2018 CeMDA.130. .. 22A。doi:10.1007/s10569-017-9805-5。
\item 珍珠，J.C.；Conrath，b.j.; 等人 (1991年)。“根据旅行者号数据确定的海王星的反照率、有效温度和能量平衡”。地球物理研究杂志。96: 18,921-930。Bibcode:1991JGR....9618921P。doi:10.1029/91JA01087。
\item 马拉玛，安东尼; 克鲁布塞克，布鲁斯; 巴夫洛夫，赫里斯托 (2017)。“行星的综合宽带震级和反照率，适用于系外行星和第九行星”。伊卡洛斯。282:19-33。arXiv:1609.05048。Bibcode:2017Icar .. 282...19米。doi:10.1016/j.icarus.2016.09.023。S2CID119307693。  
\item Mallama, A.;希尔顿，J.L.(2018)。“计算行星的表观星等天文年鉴”。天文学与计算。25:10-24.arXiv:1808.01973。Bibcode:2018A & C .... 25...10米。doi:10.1016/j.ascom.2018.08.002。S2CID69912809。  
\item “百科全书-最聪明的身体”。IMCCE。已存档从2023年7月24日开始。已检索5月29日2023。
\item Espenak，Fred (2005年7月20日)。“十二年行星星历: 1995-2006”。NASA.已存档从2012年12月5日的原件。已检索3月1日2008。
\item "海王星"。NASA科学。2017年11月10日。已检索7月19日2024。
\item Chang，Kenneth (2014年10月18日)。“我们对海王星的认识中的黑点”。纽约时报。已存档从原件到2014年10月28日。已检索10月21日2014。
\item “探索 | 海王星”。NASA太阳系探索。2017年11月10日。已存档从原件到2020年7月17日。已检索2月3日2020。1989年，NASA的旅行者2号成为第一个也是唯一一个近距离研究海王星的航天器。
\item 波多拉克，M.；魏兹曼，A.；马利，M。(1995年12月)。“天王星和海王星的比较模型”。行星与空间科学。43(12):1517-1522。Bibcode:1995P & SS。43.1517P。doi:10.1016/0032-0633(95)00061-5。
\item 鲁尼，乔纳森我。(1993年9月)。“天王星和海王星的大气层”。天文学和天体物理学年度评论。31:217-263。Bibcode:1993ARA & A .. 31 .. 217L。doi:10.1146/annurev.aa.31.090193.001245。
\item Munsell，Kirk; Smith，Harman; Harvey，Samantha (2007年11月13日)。"Neptune概述"。太阳系探测。NASA.存档自原来的2008年3月3日。已检索2月20日2008。
\item “双子座北望远镜有助于解释为什么天王星和海王星是不同的颜色-来自双子座天文台的观测，这是美国国家科学基金会NOIRLab的一个项目，其他望远镜显示，天王星上的过度薄雾使其比海王星苍白。。noirlab.edu。2022年5月31日。已存档从2022年7月30日开始。已检索7月30日2022年。
\item 香农·斯特隆(2020年12月22日)。“海王星的怪异暗斑变得更加怪异-在观察行星的大型漆黑风暴时，天文学家发现了一个较小的涡旋，他们将其命名为Dark Spot Jr”。纽约时报。已存档从原件到2020年12月22日。已检索12月22日2020。
\item Suomi, V.E.;Limaye, S.美国；约翰逊，D。R。(1991年2月22日)。“海王星的狂风: 一种可能的机制”。科学。251(4996):929-932。Bibcode:1991Sci。251。929s。doi:10.1126/science.251.4996.929。ISSN0036-8075。PMID17847386。S2CID46419483。   
\item 哈伯德，W.B.(1997)。“海王星的深层化学”。科学。275(5304):1279-80。doi:10.1126/science.275.5304.1279。PMID9064785。S2CID36248590。  
\item 内特尔曼，N.；法语，M；霍尔斯特，B.；雷德默，R。“木星，土星和海王星的内部模型”(PDF)。罗斯托克大学。存档自原来的(PDF)2011年7月18日。已检索2月25日2008。
\item 威尔福德，约翰N.(1982年6月10日)。"数据显示2个环环绕海王星"。纽约时报。已存档从2008年12月10日的原件开始。已检索2月29日2008。
\item 艾伦·赫希菲尔德 (2001)。视差: 测量宇宙的竞赛。纽约，纽约: 亨利·霍尔特。ISBN  978-0-8050-7133-7。
\item 马克·利特曼 (2004)。超越行星: 发现外太阳系。Mineola，纽约:多佛出版物。ISBN 978-0-486-43602-9。
\item 布里特，罗伯特·罗伊 (2009)。伽利略发现了海王星，新理论声称。NBC新闻。已存档从2013年11月4日的原件。已检索7月10日2009年。
\item 布瓦德，A.(1821)。表法国经度局(法语)。巴黎: Bachelier。
\item 艾里，G.B.(1846年11月13日)。“历史上与天王星外部行星的发现有关的一些情况的说明”。皇家天文学会月刊。7(10):121-144.Bibcode:1846MNRAS...7..121A。doi:10.1093/mnras/7.9.121。
\item 约翰·J·奥康纳；罗伯逊，埃德蒙·F.(2006)。“约翰·库奇·亚当斯对海王星发现的描述”。圣安德鲁斯大学。已存档从2008年1月26日开始。已检索2月18日2008。
\item 亚当斯，J.C.(1846年11月13日)。“对观察到的天王星运动不规则性的解释，关于更远的行星干扰的假设”。皇家天文学会月刊。7(9):149-152。Bibcode:1846MNRAS...7..149A。doi:10.1093/mnras/7.9.149。已存档(PDF)从2019年5月2日的原件。已检索8月25日2019。
\item Challis, J.(1846年11月13日)。“剑桥天文台探测天王星外部行星的观测记录”。皇家天文学会月刊。7(9):145-149。Bibcode:1846MNRAS...7..145C。doi:10.1093/mnras/7.9.145。已存档(PDF)从2019年5月4日开始。已检索8月25日2019。
\item 麻袋，哈拉尔德 (2017年12月12日)。“詹姆斯·查利斯和他未能发现海王星”。scihi.org。已存档从原件到2021年11月15日。已检索11月15日2021。
\item 加勒，J.G.(1846年11月13日)。“关于在柏林发现Le Verrier行星的记录”。皇家天文学会月刊。7(9): 153。Bibcode:1846MNRAS...7..153G。doi:10.1093/mnras/7.9.153。
\item Gaherty，Geoff (2011年7月12日)。“海王星完成了自1846年发现以来的第一个轨道”。space.com。已存档从原始的2019年8月25日。已检索9月3日2019。
\item 托马斯·莱文森 (2015)。寻找瓦肯人...阿尔伯特·爱因斯坦如何摧毁一颗行星，发现相对论，并破译宇宙。兰登书屋。第38页。
\item 威廉姆斯，马特 (2015年9月14日)。“气体 (和冰) 巨型海王星”。今天的宇宙。存档自原来的2023年9月27日。已检索4月26日2024。
\item Kollerstrom，Nick (2001)。“海王星的发现。英国的共同预测案例“。伦敦大学学院。存档自原来的2005年11月11日。已检索3月19日2007。
\item 威廉·希恩; 尼古拉斯·科勒斯特伦; 瓦夫，克雷格·B。(2004年12月)。“被盗行星的情况-英国人偷了海王星吗？”。科学美国人。JSTOR26060804。已存档从2011年3月19日的原件。已检索1月20日2011年。 
\item 摩尔(2000):206
\item 马克·利特曼 (2004)。超越行星，探索外太阳系。Mineola，纽约:多佛出版物。第50页。ISBN 978-0-486-43602-9。
\item 鲍姆，理查德; 希恩，威廉 (2003)。寻找瓦肯星球: 牛顿发条宇宙中的幽灵。基础书籍.pp. 109-10。ISBN 978-0-7382-0889-3。
\item 金格里奇，欧文 (1958年10月)。“天王星和海王星的命名”。太平洋天文学会传单。8(352):9-15。Bibcode:1958ASPL....8....9G。
\item 辛德，J。R。(1847)。“剑桥天文台有关新行星 (海王星) 的第二次会议报告”。天文学Nachrichten。25(21):309-314.Bibcode:1847an…… 25.309。。doi:10.1002/asna.18470252102。已存档从原件到2021年9月29日。已检索6月12日2019。
\item 福雷，冈特; 梅辛，特蕾莎·M。(2007)。“海王星: 更多惊喜”。行星科学导论。多德雷赫特: 斯普林格。pp. 385-399。doi:10.1007/978-1-4020-5544-7_19。ISBN 978-1-4020-5544-7。
\item “行星和卫星的名称和发现者”。行星命名法。美国地质调查局。2008年12月17日。已存档从原件到2018年8月9日。已检索3月26日2012。
\item “行星语言学”。nineplanets.org。已存档从2010年4月7日的原件。已检索4月8日2010。
\item "Sao Hải Vương - 'Cục băng' khổng lồ xa tít tắp"。Kenh14(越南语)。2010年10月31日。已存档从原件到2018年7月30日。已检索7月30日2018。
\item “行星的希腊名字”。2010年4月25日。已存档从2010年5月9日的原件。已检索7月14日2012。海王星或波塞冬正如它的希腊名字一样，是海洋之神。它是来自太阳的八颗行星...
\item 艾廷格，Yair (2009年12月31日)。“天王星和海王星终于得到了希伯来名字”。《国土报》。已存档从原件到2018年6月25日。已检索8月16日2018。
\item Belizovsky，Avi (2009年12月31日)。"אוראנוס הוא מהיום אורון ונפטון מעתה רהב"[天王星现在是奥伦，海王星现在是拉哈夫]。Hayadan(希伯来语)。已存档从原件到2018年6月24日。已检索8月16日2018。
\item 伊恩 (2019年9月25日)。“行星语言学 | 拉丁语，希腊语，梵语和不同的语言”。九大行星。已存档从原始的2019年2月2日。已检索2月1日2024。
\item 穆罕默德，卡迪尔 (1975)。"Waruna"。Kamus Kebangsaan Ejaan Baru, Inggeris-bahasa Malaysia, Bahasa Malaysia-Inggeris。蒂蒂旺萨。第299、857页。已存档从原件到2021年9月29日。已检索5月29日2021。
\item "Neptun"。Kamus Dewan(第四版)。Dewan Bahasa dan Pustaka马来西亚。2017。已存档从最初的2021年5月7日。已检索5月5日2021。
\item "Neptunus"。Kamus Besar Bahasa印度尼西亚(第三版)。Badan Pengembangan dan Pembinaan Bahasa Indonesia.2016。已存档从原件到2021年9月29日。已检索5月5日2021。
\item 世纪词典(1914年)
\item Long，Tony (2008年1月21日)。1979年1月21日: 海王星在冥王星古怪的轨道外移动。有线。已存档从2008年3月27日的原件开始。已检索3月13日2008。
\item "Neptune数据表"。Space.com。2006年11月17日。已存档从2023年2月11日开始。已检索2月11日2023。
\item 斯特恩，艾伦; 托伦，大卫·詹姆斯 (1997)。冥王星和Charon。亚利桑那大学出版社。pp. 206-208。ISBN 978-0-8165-1840-1。
\item 韦斯曼，保罗·R。(1995年9月)。“柯伊伯带”。天文学和天体物理学年度评论。33(1):327-357。Bibcode:1995ARA & A .. 33 .. 327W。doi:10.1146/annurev.aa.33.090195.001551。ISSN0066-4146。 
\item “冥王星的状态: 澄清”。国际天文学联合会,新闻稿。1999。存档自原来的2006年6月15日。已检索5月25日2006。
\item “天文学联盟2006年大会: 第5和第6号决议”(PDF)。IAU.2006年8月24日。已存档(PDF)从2008年6月25日的原件。已检索7月22日2008。
\item Uns ö ld，Albrecht; Baschek，Bodo (2001年)。新宇宙: 天文学和天体物理学导论(第5版)。斯普林格.表3.1，第47页。Bibcode:2001年ncia。书 ..... U。ISBN  978-3-540-67877-9。
\item 老板，艾伦·P。(2002)。“气体和冰巨行星的形成”。地球和行星科学快报。202(3-4):513-23.Bibcode:2002E和psl.202.513B。doi:10.1016/S0012-821X(02)00808-7。
\item 洛维斯，C.；市长，M.；阿里伯特，Y.；奔驰，W.(2006年5月18日)。“海王星三重奏和他们的腰带”。ESO。已存档从2010年1月13日的原件。已检索2月25日2008。
\item Atreya, S.;埃格勒，P.；贝恩斯，K。(2006)。“天王星和海王星上的水-氨离子海洋？”(PDF)。地球物理研究文摘。8。05179。已存档(PDF)从2012年2月5日的原件。已检索11月7日2007。
\item Shiga，David (2010年9月1日)。“奇怪的水潜伏在巨大的行星”。新科学家。2776号.已存档从原件到2018年2月12日。已检索2月11日2018。
\item 克尔，理查德·A.(1999年10月)。“海王星可能会将甲烷粉碎成钻石”。科学。286(5437): 25a-25。doi:10.1126/science.286.5437.25a。PMID10532884。S2CID42814647。  
\item 莎拉·卡普兰 (2017年8月25日)。“它在天王星和海王星上下雨固体钻石”。华盛顿邮报。已存档从原始的2017年8月27日。已检索8月27日2017。
\item 克劳斯，D.; 等人 (2017年9月)。“行星内部条件下激光压缩碳氢化合物中钻石的形成”(PDF)。自然天文学。1(9):606-11.Bibcode:2017NatAs...1 .. 606K。doi:10.1038/s41550-017-0219-9。OSTI1393331。S2CID46945778。已存档(PDF)从原始的2018年7月25日。已检索8月25日2018。   
\item 凯恩，肖恩 (2016年4月29日)。“闪电风暴使土星和木星上的钻石下雨”。商业内幕。已存档从原件到2019年6月26日。已检索5月22日2019。
\item 鲍德温，艾米丽 (2010年1月21日)。“钻石海洋可能在天王星和海王星上”。天文学现在。存档自原来的2013年12月3日
\item 布拉德利，D.K.；艾格特，J.H.；希克斯，d.g.；Celliers，下午.(2004年7月30日)。“将金刚石压缩成导电流体的冲击”(PDF)。物理评论信。93(19) 195506。Bibcode:2004PhRvL .. 93s5506B。doi:10.1103/physrevlett.93.195506。高密度脂蛋白:1959.3/380076。PMID15600850。S2CID6203103。存档自原来的(PDF)2016年12月21日。已检索3月16日2016。   
\item 艾格特，J.H.；希克斯，d.g.；Celliers，下午；Bradley，D.K.; 等人 (2009年11月8日)。“超高压下金刚石的熔化温度”。自然物理。6(40):40-43。Bibcode:2010NatPh...6...40E。doi:10.1038/nphys1438。
\item 安德鲁斯，罗宾·乔治 (2023年8月18日)。海王星的云层已经消失，科学家认为他们知道为什么-最近的一项研究表明太阳周期与太阳系第八颗行星的大气之间存在关系。纽约时报。已存档从原始的2023年8月18日。已检索8月21日2023。
\item 酥脆，D.；哈默尔，H.B.(1995年6月14日)。“哈勃太空望远镜观测海王星”。哈勃新闻中心。已存档从2007年8月2日的原件。已检索4月22日2007。
\item 贝基·费雷拉 (2024年1月4日)。“天王星和海王星揭示了它们的真实颜色-海王星并不像你所相信的那样蓝，在新的研究中，天王星的颜色变化得到了更好的解释”。纽约时报。已存档从2024年1月5日开始。已检索1月5日2024。
\item NASA科学编辑团队 (2022年5月31日)。“为什么天王星和海王星是不同的颜色”。NASA.已存档从原来的2023年10月31日。已检索10月30日2023。
\item 威廉姆斯，马特 (2015年9月14日)。“气体 (和冰) 巨型海王星”。Phys.org。已检索8月27日2024。
\item 埃尔金斯-坦顿，琳达T.(2006)。天王星、海王星、冥王星和外太阳系。纽约: 切尔西之家。pp。 79-83。ISBN 978-0-8160-5197-7。
\item 麦克斯，C.E.;麦金塔，B。A.;吉巴德，S.G.;木槌，D.T.; 等人 (2003年1月)。“用凯克望远镜自适应光学观测到的海王星上的云结构”。天文杂志。125(1):364-375。Bibcode:2003AJ....125..364M。doi:10.1086/344943。
\item Gianopoulos，Andrea (2023年8月16日)。海王星消失的云与太阳周期有关。NASA。已存档从2023年8月24日开始。已检索8月24日2023。
\item 查韦斯，埃兰迪; 德佩特，伊姆克; 雷德温，艾琳; 莫尔特，爱德华M.；罗曼，迈克尔·T；佐尔齐，安德里亚; 阿尔瓦雷斯，卡洛斯; 坎贝尔，兰迪; 德克莱尔，凯瑟琳; 韦索，里卡多; 黄，迈克尔·H；盖茨，埃莉诺; 林纳姆，保罗·大卫; 戴维斯，阿什利·G。；艾科克，乔尔; 麦克罗伊，杰森; 佩尔蒂埃，约翰; 里德诺，安东尼; 斯蒂克，特里 (2023年11月1日)。“从1994年到2022年海王星在近红外波长的演变”。伊卡洛斯。404115667。arXiv:2307.08157。Bibcode:2023年Icar .. 40415667C。doi:10.1016/j.icarus.2023.115667。S2CID259515455。已存档从2023年8月24日开始。已检索8月24日2023。我们发现云活动与太阳lyman-alpha (121.56 nm) 之间存在明显的正相关辐照度为海王星云活动的周期性是由太阳紫外线发射触发的光化学云/霾产生引起的理论提供了支持。 
\item 恩克雷纳兹，泰雷泽 (2003年2月)。ISO对巨行星和土卫六的观测: 我们学到了什么？行星与空间科学。51(2):89-103.Bibcode:2003P & SS...51...89E。doi:10.1016/S0032-0633(02)00145-9。
\item Broadfoot, A.L.;Atreya, S.K.;Bertaux，J.L.; 等人 (1999年)。“海王星和海卫一的紫外光谱仪观测”(PDF)。科学。246(4936):1459-66.Bibcode:1989Sci。246。1459B。doi:10.1126/science.246.4936.1459。PMID17756000。S2CID21809358。已存档(PDF)从2008年5月28日的原件。已检索3月12日2008。   
\item 赫伯特，弗洛伊德; 桑德尔，比尔·R。(1999年8月至9月)。“天王星和海王星的紫外线观测”。行星与空间科学。47(8-9): 1,119-139。Bibcode:1999P & SS。47.1119h。doi:10.1016/S0032-0633(98)00142-1。
\item 欧文，帕特里克G.J.;多宾森，杰克; 詹姆斯，阿朱那; 提比，尼古拉斯·A；艾米·西蒙；Fletcher, Leigh N.;罗曼，迈克尔·T；格伦·S·奥顿；黄，迈克尔·H; 托莱多，丹尼尔; 佩雷斯-霍约斯，圣地亚哥; 贝克，朱莉 (2024年2月)。“模拟天王星的颜色和大小的季节性周期，并与海王星进行比较”。皇家天文学会月刊。527(4):11521-11538。doi:10.1093/mnras/stad3761。高密度脂蛋白:20.500.11850/657542。
\item "PIA01492的目录页"。photojournal.jpl.nasa.gov。已存档从2023年7月22日开始。已检索2月5日2024。
\item “天王星和海王星之间的细微色差”。行星社会。已存档从2024年2月5日开始。已检索2月5日2024。
\item 牛津大学。“新图像揭示了海王星和天王星的真实面貌”。phys.org。已存档从2024年2月5日开始。已检索2月5日2024。
\item 斯坦利，萨宾；杰里米·布洛瑟姆 (2004年3月)。“对流区几何形状是天王星和海王星异常磁场的原因”。自然。428(6979):151-153.Bibcode:2004 Natur.428..151S。doi:10.1038/nature02376。ISSN0028-0836。PMID15014493。S2CID33352017。   
\item 康纳尼，J.E.P.；阿库尼亚，马里奥·H; 内斯，诺曼·F。(1991)。“海王星的磁场”。地球物理研究杂志。96: 19,023-42。Bibcode:1991JGR....9619023C。doi:10.1029/91JA01165。
\item 内斯，诺曼·F；马里奥·H·阿库尼亚; 伦纳德·F·伯拉加；康纳尼，约翰E。P.;罗纳德·P·莱普；纽鲍尔，弗里茨·M。(1989年12月15日)。“海王星的磁场”。科学。246(4936):1473-1478。Bibcode:1989Sci。246.1473n。doi:10.1126/science.246.4936.1473。ISSN0036-8075。PMID17756002。S2CID20274953。已存档(PDF)从原始的2019年7月10日。已检索8月25日2019。   
\item 罗素，C.T.；卢曼，J.G.(1997)。“海王星: 磁场和磁层”。加州大学洛杉矶分校。已存档从原件到2019年6月29日。已检索8月10日2006。
\item 拉米，L。(2020年11月9日)。“天王星和海王星的极光排放”。皇家学会的哲学交易A: 数学，物理和工程科学。378(2187) 20190481。皇家学会。Bibcode:2020RSPTA.37890481L。doi:10.1098/rsta.2019.0481。PMC7658782。PMID33161867。  
\item “ESA门户-火星快车在火星上发现极光”。欧洲航天局。2004年8月11日。已存档从原件到2012年10月19日。已检索8月5日2010。
\item “NASA的韦伯首次捕捉到海王星的极光”。NASA韦伯任务小组.2025年3月26日。已检索3月26日2025。
\item Lavoie，Sue (1998年1月8日)。"PIA01142: 海王星滑板车"。NASA.已存档从原来的2013年10月29日。已检索3月26日2006。
\item 哈默尔，H. B.；Beebe, R.F.;De Jong, E.M、；汉森，C。等 (1989年9月22日)。“通过跟踪旅行者图像中的云层获得的海王星风速”。科学。245(4924):1367-1369。Bibcode:1989Sci...245.1367H。doi:10.1126/science.245.4924.1367。ISSN0036-8075。PMID17798743。S2CID206573894。   
\item 伯吉斯(1991):64-70。
\item 雷·维拉德; 特里·德维特 (2003年5月15日)。“更明亮的海王星暗示着季节的行星变化”。哈勃新闻中心。已存档从2008年2月28日的原件。已检索2月26日2008。
\item Lavoie，Sue (2000年2月16日)。“PIA02245: 海王星的蓝绿色大气”。NASA JPL。已存档从原来的2013年8月5日。已检索2月28日2008。
\item 奥顿，G.美国；Encrenaz, T.；莱拉特，C.；Puetter，R.; 等人 (2007年10月)。“甲烷逸出以及海王星大气温度强烈的季节性和动力扰动的证据”(PDF)。天文学与天体物理学。473(1):L5-L8。Bibcode:2007A & A...473L...5O。doi:10.1051/0004-6361:20078277。ISSN0004-6361。S2CID54996279。已存档(PDF)从2024年2月1日起。已检索2月1日2024。   
\item 格伦奥顿；泰雷斯·恩克雷纳兹(2007年9月18日)。温暖的南极？是的，在海王星上!”。ESO。已存档从2010年3月23日的原件。已检索9月20日2007。
\item 哈默尔，H. B.；洛克伍德，G。W.;米尔斯，J。R、；巴内特，C。D.(1995)。“1994年哈勃太空望远镜对海王星云结构的成像”。科学。268(5218):1740-42。Bibcode:1995Sci。268.1740h。doi:10.1126/science.268.5218.1740。PMID17834994。S2CID11688794。  
\item Lavoie，Sue (1996年1月29日)。“PIA00064: 海王星的黑点 (D2) 在高分辨率”。NASA JPL。已存档从原来的2013年9月27日。已检索2月28日2008。
\item 神秘的海王星黑斑首次从地球上探测到。ESO。已存档从2023年8月26日开始。已检索8月26日2023。
\item SDS-2015这意味着它是2015年发现的南部黑斑。H.黄，迈克尔; 托勒夫森，约书亚; 我。徐，安德鲁; 德佩特，伊姆克; A。西蒙，艾米; 韦索，里卡多; 桑切斯-拉维加，奥古斯丁; 斯罗莫夫斯基，劳伦斯; 弗莱，帕特里克; Luszcz-cook，Statia (2018年2月15日)。“海王星上的新黑暗漩涡”。美国天文学会。155(3)。摘要部分。doi:10.3847/1538-3881(2025年7月1日不活跃)。
\item S.G.,吉巴德; 德佩特，I.；罗伊，h.g.；马丁，S.; 等人 (2003年)。“来自高空间分辨率近红外光谱的海王星云特征的高度”(PDF)。伊卡洛斯。166(2):359-74.Bibcode:2003年Icar。。166。359克。doi:10.1016/j.icarus.2003.07.006。存档自原来的(PDF)2012年2月20日。已检索2月26日2008。
\item 斯特拉特曼，P.W.；Showman, A.P.;道林，t.e.；Sromovsky, L.A.(2001)。“史诗般的模拟明亮的同伴到海王星的大黑点”(PDF)。伊卡洛斯。151(2):275-85。Bibcode:1998年伊卡。。132。239L。doi:10.1006/icar.1998.5918。已存档(PDF)从2008年2月27日的原件开始。已检索2月26日2008。
\item Sromovsky, L.A.;弗莱，下午；道林，t.e.；贝恩斯，K.H.(2000)。“海王星上新黑点的异常动态”。美国天文学会公报。32: 1005。Bibcode:2000DPS....32.0903S。
\item "生日快乐海王星"。欧空局/哈勃。已存档从2011年7月15日开始。已检索7月13日2011年。
\item 博鲁基，W.J.(1989)。“海王星闪电活动的预测”。地球物理研究快报。16(8):937-939。Bibcode:1989年乔治。。16 .. 937B。doi:10.1029/gl016i008p00937。
\item 阿普林，K.L.；费舍尔，G.；诺德海姆，T.A.；科诺瓦连科，A.；Zakharenko, V.;Zarka，P。(2020)。“冰巨人的大气电”。空间科学评论。216(2): 26。arXiv:1907.07151。Bibcode:2020年SSRv。。216。26A。doi:10.1007/s11214-020-00647-0。
\item Gurnett, D.A.;Kurth, W.美国；凯恩斯，我。H、; 格兰罗斯，L。J.(1990)。“海王星磁层中的哨声: 大气闪电的证据”。地球物理研究杂志: 空间物理。95(A12):20967-20976。Bibcode:1990JGR....9520967G。doi:10.1029/ja095ia12p20967。高密度脂蛋白:2060/19910002329。
\item 贝尔彻，j.w.；Bridge, H.S.;巴格纳尔，F.；科皮，B.；潜水员，O.；埃维塔尔，A.；戈登，G.S.；拉撒路，A.J.；麦克纳特，R.L.；奥格尔维，k.w.；理查森，J.D.；西斯科，g.l.；西特勒，E.C.；斯坦伯格，J.T.；沙利文，J.D.；萨博，A.；维拉纽瓦，L.；Vasyliunas, V.M.;张，M。(1989)。“海王星附近的等离子体观测: 旅行者2号的初步结果”。科学。246(4936):1478-1483。Bibcode:1989Sci。246.1478B。doi:10.1126/science.246.4936.1478。PMID17756003。 
\item 吉巴德，s.g.；利维，E.H.；鲁宁，J.I.; 德佩特，I.(1999)。“海王星上的闪电”。伊卡洛斯。139(2):227-234。Bibcode:1999年Icar。。139。227克。doi:10.1006/icar.1999.6101。
\item Aglyamov, Y.S.;鲁尼，J.；Atreya, S.;Guillot, T.;贝克尔，H.N.；莱文，S.；博尔顿，s.j.(2023年6月)。“非理想气体中的巨型行星闪电”。行星科学杂志。4(6) 111。Bibcode:2023PSJ.....4..111A。doi:10.3847/PSJ/acd750。
\item Lindal, Gunnar F.(1992)。“海王星的大气-对旅行者2号获得的无线电掩星数据的分析”。天文杂志。103:967-82。Bibcode:1992AJ....103..967L。doi:10.1086/116119。
\item “第12类-巨型行星-热量和形成”。3750-行星，卫星和环。科罗拉多大学，博尔德。2004。已存档从原件到2008年6月21日。已检索3月13日2008。
\item 佩特，伊姆克·德; 利绍尔，杰克·J。(2001年12月6日)。行星科学。剑桥大学出版社，第224页。ISBN 978-0-521-48219-6。已存档从原件到2021年9月29日。已检索3月15日2023。
\item Meeus，Jean (1998)。天文算法。弗吉尼亚州里士满: 威尔曼-贝尔。第273页。通过进一步使用vsop87来补充。
\item "行星概况"。已存档从原始的2024年2月2日。已检索1月2日2024。
\item 麦基，罗宾 (2011年7月9日)。“海王星的第一个轨道: 天文学的转折点”。守护者。已存档从2016年8月23日起。已检索12月15日2016。
\item 南希阿特金森 (2010年8月26日)。“清除海王星轨道上的混乱”。今天的宇宙。已存档从2023年9月29日开始。已检索2月1日2024。
\item Lakdawalla，Emily [@ elakdawalla] (2010年8月19日)。"Doh!RT @ lukedones: 来自JPL的Bill Folkner: 海王星将达到与1846年9月23日相同的黄道经度，即2011年7月12日。(鸣叫)-通过Twitter。
\item "海王星2010-2011的地平线输出"。2007年11月16日。存档自原来的2013年5月2日。已检索2月25日2008。-使用太阳系动力学组生成的数字，地平线在线星历系统。
\item 威廉姆斯，大卫R.(2005年1月6日)。"行星概况"。NASA.已存档从2008年9月25日开始。已检索2月28日2008。
\item 哈伯德，W。B.;Nellis, W.J.;米切尔，A.C.;福尔摩斯，N.C.; 等人 (1991年8月9日)。“海王星的内部结构: 与天王星的比较”。科学。253(5020):648-651。Bibcode:1991Sci。253。648小时。doi:10.1126/science.253.5020.648。ISSN0036-8075。PMID17772369。S2CID20752830。已存档从2018年10月23日的原件。已检索6月12日2019。   
\item Thommes, E.W.;邓肯，M。J.;利维森，h.f.(2002年5月)。“木星和土星之间天王星和海王星的形成”。天文杂志。123(5):2862-2883。arXiv:astro-ph/0111290。Bibcode:2002AJ....123.2862T。doi:10.1086/339975。S2CID17510705。 
\item 凯瑟琳·汉森 (2005年6月7日)。“早期太阳系的轨道洗牌”。地理时间。已存档从2007年9月27日的原件。已检索8月26日2007。
\item 克里达，A.(2009)。“太阳系形成”。现代天文学评论。第二十一卷。第3008页。arXiv:0903.3008。Bibcode:2009RvMA...21 .. 215C。doi:10.1002/9783527629190.ch12。ISBN  978-3-527-62919-0。S2CID 118414100。
\item 您将收到Desch，s.j.(2007)。“太阳星云中的质量分布和行星形成”(PDF)。天体物理学杂志。671(1):878-93。Bibcode:2007ApJ...671 .. 878D。doi:10.1086/522825。S2CID120903003。存档自原来的(PDF)2020年2月7日。  
\item 史密斯，R.；丘吉尔，L。J.;怀亚特，M。C.;Moerchen，M。M.; 等人 (2009年1月)。“解决了 η Telescopii周围的碎片圆盘发射: 一个年轻的太阳系或正在进行的行星形成？”天文学与天体物理学。493(1):299-308。arXiv:0810.5087。Bibcode:2009A & A...493 .. 299S。doi:10.1051/0004-6361:200810706。ISSN0004-6361。S2CID6588381。  
\item 斯特恩，S.艾伦; 科尔威尔，约书亚·E。(1997)。“原始埃奇沃思-柯伊伯带的碰撞侵蚀和30-50 AU柯伊伯带的产生”。天体物理学杂志。490(2):879-82。Bibcode:1997ApJ。490 .. 879S。doi:10.1086/304912。
\item 佩蒂特，让-马克; 莫比德利，亚历山德罗; 瓦尔塞奇，乔瓦尼B。(1999)。“大的分散小行星和小身体带的激发”(PDF)。伊卡洛斯。141(2):367-87.Bibcode:1999年Icar .. 141 .. 367页。doi:10.1006/icar.1999.6166。存档自原来的(PDF)2007年12月1日。已检索6月23日2007。
\item "Transneptunian对象列表"。小行星中心。已存档从2010年10月27日的原件。已检索10月25日2010。
\item Jewitt，David (2004)。《Plutinos》。加州大学洛杉矶分校。已存档从2007年4月19日的原件。已检索2月28日2008。
\item 瓦拉迪，F.(1999)。“3:2轨道共振中的周期轨道及其稳定性”。天文杂志。118(5):2526-31.Bibcode:1999AJ....118.2526V。doi:10.1086/301088。
\item 约翰·戴维斯 (2001)。超越冥王星: 探索太阳系的外部界限。剑桥:剑桥大学出版社。p。 104。ISBN 978-0-521-80019-8。
\item 蒋，E。I.;约旦，A。B.;米利斯，R。L.;Buie, M.W.; 等人 (2003年7月)。“柯伊伯带的共振占领: 5:2和特洛伊共振的案例”。天文杂志。126(1):430-443。arXiv:astro-ph/0301458。Bibcode:2003AJ....126..430C。doi:10.1086/375207。ISSN0004-6256。S2CID54079935。  
\item 谢泼德，斯科特·S.；特鲁希略，查德威克A。(2010年9月10日)。"检测到尾随 (L5) Neptune特洛伊木马"。科学。329(5997): 1304。Bibcode:2010Sci...329.1304S。doi:10.1126/science.1189666。ISSN0036-8075。PMID20705814。S2CID7657932。   
\item De La Fuente Marcos, C.&De La Fuente Marcos, R.(2012)。“(309239) 2007 RW10: 海王星的大型临时准卫星”。天文学和天体物理学快报。545(2012): l9。arXiv:1209.1577。Bibcode:2012A & A...545L...9D。doi:10.1051/0004-6361/201219931。S2CID118374080。 
\item “新天王星和海王星卫星”。地球与行星实验室。卡内基科学研究所。2024年2月23日已存档从原始的2024年2月23日。已检索2月23日2024。
\item 克雷格·B·阿格诺；汉密尔顿，道格拉斯P.(2006年5月)。“海王星在双行星引力相遇中捕获了它的月亮Triton”。自然。441(7090):192-194.Bibcode:2006 Natur.441..192A。doi:10.1038/nature04792。ISSN0028-0836。PMID16688170。S2CID4420518。   
\item Chyba, Christopher F.;扬科夫斯基，d.g.；尼科尔森，P.D.(1989)。“海王星-海王星系统的潮汐演化”。天文学和天体物理学。219(1-2):L23-L26。Bibcode:1989A & A...219L..23C。
\item 威尔福德，约翰N.(1989年8月29日)。“Triton可能是太阳系中最冷的点”。纽约时报。已存档从2008年12月10日的原件。已检索2月29日2008。
\item “Triton-nasa科学”。2017年11月21日。已存档从2024年1月7日开始。已检索1月7日2024。
\item 罗伯特·M·纳尔逊；史密斯，威廉·D；沃利斯，布拉德·D；Horn，Linda J.; 等人 (1990年10月19日)。“海王星卫星Triton表面的温度和热发射率”。科学。250(4979):429-431。Bibcode:1990科学。250 .. 429N。doi:10.1126/science.250.4979.429。ISSN0036-8075。PMID17793020。S2CID20022185。   
\item “12.3: 泰坦和特里顿”。2016年10月7日。已存档从2024年1月7日开始。已检索1月7日2024。
\item “海王星: 海王星的月亮”。2010年1月。已存档从2024年1月7日开始。已检索1月7日2024。
\item 斯通，E.C.;矿工，E。D.(1989年11月15日)。“旅行者2号与海王星系统相遇”。科学。246(4936):1417-1421。Bibcode:1989Sci。246.1417s。doi:10.1126/science.246.4936.1417。ISSN0036-8075。PMID17755996。S2CID9367553。   
\item 布朗，迈克尔E。“矮行星”。加州理工学院地质科学系。存档自原来的2011年7月19日。已检索2月9日2008。
\item 霍尔曼，m.j.;Kavelaars, J.J.；Grav，T.; 等人 (2004年)。“发现海王星的五颗不规则卫星”(PDF)。自然。430(7002):865-67.Bibcode:2004 Natur.430..865H。doi:10.1038/nature02832。PMID15318214。S2CID4412380。已存档(PDF)从2013年11月2日的原件。已检索10月24日2011年。   
\item “海王星的五个新卫星”。BBC新闻。2004年8月18日。已存档从2007年8月8日的原件。已检索8月6日2007。
\item Grush，Loren (2019年2月20日)。“海王星的新发现的月亮可能是古代碰撞的幸存者”。The Verge。已存档从原件到2019年2月21日。已检索2月22日2019。
\item O “卡拉汉，乔纳森 (2022年9月21日)。“海王星及其环与韦伯望远镜聚焦-来自太空天文台的新图像提供了红外行星的新视角”。纽约时报。已存档从2022年9月22日开始。已检索9月23日2022年。
\item 克鲁克申克，戴尔·P。(1996)。海王星和Triton。亚利桑那大学出版社。pp。 703-804.ISBN 978-0-8165-1525-7。
\item 蓝，詹妮弗 (2004年12月8日)。“命名环和环间隙命名”。行星命名法地名。美国地质勘探局。已存档从原来的2010年7月5日。已检索2月28日2008。
\item 桂南，e.f.；哈里斯，C.；马洛尼，f.p.(1982)。“海王星环系统的证据”。美国天文学会公报。14: 658。Bibcode:1982BAAS...14..658G。
\item Goldreich, P.;Tremaine, S.;边界，N。(1986年8月)。“走向海王星弧环的理论”(PDF)。天文杂志。92: 490。Bibcode:1986AJ.....92..490G。doi:10.1086/114178。于2021年8月9日从原件存档。已检索6月12日2019。
\item 菲利普·D·尼科尔森；库克，马伦·L；Matthews，Keith; Elias，Jonathan H.; Gilmore，Gerard; 等 (1990年9月)。“海王星的五个恒星掩星: 对环弧的进一步观察”。伊卡洛斯。87(1):1-39。Bibcode:1990年伊卡尔。87。1N。doi:10.1016/0019-1035(90)90020-A。
\item 考克斯，亚瑟·N.(2001)。艾伦的天体物理量。斯普林格.ISBN 978-0-387-98746-0。
\item Munsell，Kirk; Smith，Harman; Harvey，Samantha (2007年11月13日)。"行星: 海王星: 环"。太阳系探测。NASA.已存档从2012年7月4日开始。已检索2月29日2008。
\item 萨洛，海基; 汉尼宁，Jyrki (1998)。“海王星的部分环: 加拉太对自引力弧粒子的作用”。科学。282(5391):1102-04。Bibcode:1998Sci。282.1102s。doi:10.1126/science.282.5391.1102。PMID9804544。 
\item “海王星的光环正在消失”。新科学家。2005年3月26日。已存档从2008年12月10日的原件。已检索8月6日2007。
\item 施穆德，R.W.小; 贝克，R.E.；福克斯，J.；Krobusek, B.A.;帕夫洛夫，H.; 马拉马，A。(2016年3月29日)。海王星的世俗和旋转亮度变化 (未发表的手稿)。arXiv:1604.00518。
\item 有关震级数据，请参阅相应的文章。
\item 摩尔(2000):207。
\item 例如，在1977年，甚至海王星的旋转周期仍然不确定。克鲁克申克，d.p.(1978年3月1日)。“关于海王星的旋转周期”。天体物理学杂志快报。220:L57-L59。Bibcode:1978年apj... 220L..57C。doi:10.1086/182636。
\item 马克斯，C.；麦金塔，B.；吉巴德，S.；Roe，H.; 等人 (1999年)。“海王星和泰坦的自适应光学成像与W.M.凯克望远镜 ”。美国天文学会公报。31: 1512。Bibcode:1999AAS...195.9302M。
\item Nemiroff, R.;邦内尔，J.，编辑。(2000年2月18日)。“通过自适应光学的海王星”。当天的天文图片。NASA。
\item 罗迪尔，F.；罗迪尔，C.；Brahic, A.;大仲马，C.；格雷夫斯，J。E.;诺斯科特，M。J.;欧文，T.(1997年8月1日)。“海王星和变形杆菌的首次地面自适应光学观测”。行星与空间科学。45(8):1031-1036。Bibcode:1997P & SS。45.1031r。CiteSeerX10.1.1.66.7754。doi:10.1016/S0032-0633(97)00026-3。已存档从2024年2月1日起。已检索2月1日2024。 
\item Engvold，Oddbjorn (2007年5月10日)。2003-2005年天文学报告 (IAU XXVIA): IAU Transactions XXVIA。剑桥大学出版社。147f页。ISBN 978-0-521-85604-1。已存档从2023年5月11日开始。已检索3月15日2023。
\item 吉巴德，S.G.;罗伊，H.; 德佩特，I.；Macintosh，B.; 等人 (2002年3月)。“凯克望远镜对海王星的高分辨率红外成像”。伊卡洛斯。156(1):1-15。Bibcode:2002年Icar。。156 .... 1克。doi:10.1006/icar.2001.6766。ISSN0019-1035。已存档从2018年10月23日的原件。已检索6月12日2019。 
\item 辛西娅·菲利普斯(2003年8月5日)。“对遥远世界的迷恋”。SETI研究所。存档自原来的2007年11月3日。已检索10月3日2007。
\item 伯吉斯(1991):46-55。
\item Standage 2000,第188页
\item 克里斯·格布哈特; 杰夫·戈德德 (2011年8月20日)。发射34年后，旅行者2号继续探索。NASASpaceflight。已存档从2016年2月19日的原件。已检索1月22日2016。
\item 吴伟仁; 于，登云; 黄，江川; 宗，秋刚; 王，池; 于，国斌; 何，荣伟; 王，钱; 康，严; 孟，林芝；吴柯; 何建森; 李辉 (2019年1月9日)。“探索太阳系边界”。科学信息。49(1): 1。doi:10.1360/N112018-00273。
\item 安德鲁·琼斯 (2021年4月16日)。“中国将向太阳系边缘发射一对航天器”。SpaceNews。已存档从原件到2021年5月15日。已检索4月29日2021。
\item 克拉克，斯蒂芬 (2015年8月25日)。“天王星，海王星在NASA的新机器人任务中”。现在太空飞行。已存档从原来的2015年9月6日。已检索9月7日2015。
\item Spilker, T.R.;英格索尔，a.p.(2004)。“来自航空捕获视觉任务的海王星系统中的杰出科学”。美国天文学会公报。36: 1094。Bibcode:2004DPS....36.1412S。
\item 汉森，坎迪斯; 等人。“Argo-穿越外太阳系的航行”(PDF)。SpacePolicyOnline.com。空间和技术政策集团有限责任公司。存档自原来的(PDF)2015年9月24日。已检索8月5日2015。
\item “用三叉戟探索特里顿: 发现级任务”(PDF)。大学空间研究协会。2019年3月23日。已存档(PDF)从原件到2020年8月2日。已检索3月26日2019。
\item 莱默，阿比盖尔; 克莱德，布伦达; 鲁尼恩，柯比 (2020年8月)。“海王星奥德赛: 海王星-特里顿系统的使命”(PDF)。NASA科学。已存档(PDF)从原件到2020年12月15日。已检索4月18日2021。
\item 起源，世界和生命: 2023-2032年行星科学和天体生物学的十年战略(出版前版)。国家科学院出版社。2023年1月19日第800页。Bibcode:2023年猫头鹰。书 ..... N。doi:10.17226/26522。ISBN 978-0-309-47578-5。S2CID 248283239。已检索4月30日2022年。
\item 汉森-科哈切克，坎迪斯; 菲尔豪尔，卡尔 (2021年6月7日)。“Triton海洋世界测量师概念研究”(PDF)。NASA。已存档(PDF)从2023年11月9日开始。已检索1月11日2024。
\item 斯特克尔，阿曼达; 康拉德，杰克·威廉; 德卡斯克，杰森; 多兰，悉尼; 唐尼，布兰娜·格蕾丝; 费尔顿，瑞安; 汉森，薰衣草·艾尔; 吉斯切，阿琳娜; 霍瓦特，泰勒；麦克斯韦，瑞秋; 沙姆韦，安德鲁·奥; 西迪克，阿纳米卡; 斯特罗姆，迦勒; 提斯，布朗温；托德，杰西卡; 特林，凯文·T; 维莱兹，迈克尔·A; 沃尔特，卡勒姆·安德鲁; 劳斯，莱斯利·L; 哈德森，特洛伊; 史高丽，詹妮弗·E。C.(2023年12月12日)。“鹦鹉螺的科学案例: 对Triton的多次飞越任务概念”。AGU - Agu23(3089)。AGU。Bibcode:2023AGUFM.P23D3089S。已存档从2024年1月11日开始。已检索1月11日2024。
\item “中国外太阳系探测计划”。行星社会。2023年12月21日。已检索4月18日2024。
\end{enumerate}
\subsection{参考书目}
\begin{itemize}
\item 伯吉斯，埃里克 (1991)。遥远的相遇: 海王星系统。纽约:哥伦比亚大学出版社。ISBN 978-0-231-07412-4。
\item 帕特里克·摩尔(2000)。天文学数据手册。巴吞鲁日:泰勒和弗朗西斯集团。ISBN 978-0-7503-0620-1。
\end{itemize}
\subsection{进一步阅读}
\begin{itemize}
\item 矿工，埃利斯·D；Wessen, Randii R.(2002)。海王星: 行星，环和卫星。斯普林格-天文学和空间科学方面的实践书籍。伦敦; 纽约:斯普林格。ISBN 978-1-85233-216-7。
\item Standage，汤姆 (2000)。海王星文件: 天文学竞争和行星狩猎先驱的故事。纽约:沃克。ISBN 978-0-8027-1363-6。
\end{itemize}
\subsection{外部链接}

\begin{itemize}
\item NASA的海王星概况介绍 已存档2010年7月1日在回程机
\item 海王星Bill Arnett的nineplanets.org
\item 海王星 天文学演员第63集，包括完整的成绩单。
\item Neptune配置文件(2002年11月15日存档)NASA太阳系探索基地
\item 海王星及其内部卫星的交互式3D重力模拟 已存档2020年9月22日在回程机
\end{itemize}